\begin{pastexam}[label={px:ntu-110-comm-dsa-30}]{臺大電信所 110 學年度（B）第 30 題}
（單選；每題 2 分，答錯不倒扣）

In this exam,
\begin{itemize}
\item The height of a leaf (in either a tree or heap) is defined as 0.
\item The depth of the root (in either a tree or heap) is defined as 0.
\item The degree of a node (in either a tree or heap) is defined as the number of its children.
\item The size of a tree/heap is defined as the total number of keys (items) it stores.
\item All heaps, unless otherwise specified, are min heaps.
\end{itemize}

Consider a binary heap that results from successively inserting keys 6, 3, 7, 8, 5, 1, 2, 4, 9 into an initially empty heap. Then the depth of key 5 is

A.~0 \qquad B.~1 \qquad C.~2 \qquad D.~3 \qquad E.~None of above
\end{pastexam}

\begin{pastsoln}
答案是 C（2）。

本卷的堆積預設是最小堆積。 每次插入放在陣列最後一格，再和父節點比較，比父節點小就交換往上浮
（資結書第~8.1.5~節）。 陣列從索引 $1$ 算，節點 $i$ 的父節點是 $\lfloor i/2\rfloor$：
\begin{center}
\begin{tabular}{cll}
\toprule
插入 & 上浮過程 & 陣列 \\
\midrule
6 & — & $[6]$ \\
3 & $3 < 6$，換 & $[3, 6]$ \\
7 & $7 > 3$ & $[3, 6, 7]$ \\
8 & $8 > 6$ & $[3, 6, 7, 8]$ \\
5 & 在第 5 格，$5 < 6$ 換到第 2 格；$5 > 3$ 停 & $[3, 5, 7, 8, 6]$ \\
1 & 第 6 格，$1 < 7$ 換到第 3 格；$1 < 3$ 換到第 1 格 & $[1, 5, 3, 8, 6, 7]$ \\
2 & 第 7 格，$2 < 3$ 換到第 3 格；$2 > 1$ 停 & $[1, 5, 2, 8, 6, 7, 3]$ \\
4 & 第 8 格，$4 < 8$ 換到第 4 格；$4 < 5$ 換到第 2 格；$4 > 1$ 停 & $[1, 4, 2, 5, 6, 7, 3, 8]$ \\
9 & 第 9 格，$9 > 5$ & $[1, 4, 2, 5, 6, 7, 3, 8, 9]$ \\
\bottomrule
\end{tabular}
\end{center}
最後 $5$ 在第 $4$ 格。 第 $i$ 格的深度是 $\lfloor\lg i\rfloor$，第 $4$ 格深度 $2$（根 $1$ 的子節點 $4$ 的子節點）。
選 C。 注意 $5$ 原本在深度 $1$（第 $2$ 格），是插入 $4$ 時被換下去的。
\end{pastsoln}
